\begin{statement}[C.2.1; Proved: exact nomes and continuation radius]
For $w=e^{-\pi ic}$, the fixed branches give
\begin{align*}
 \kappa_1&=c+\tfrac12-\frac{3i}{2\pi}\log48,&
 q_1&=-w^2/48^3,\\
 \kappa_2&=c+\tfrac32-\frac{3i}{\pi}\log6,&
 q_2&=iw/216,\qquad q_1=(27/64)q_2^2.
\end{align*}
The family is parameterized by $0<|q_2|<1$ and exists precisely in
that domain. Writing $J(q)=q^{-1}+744+\cdots$, one has
\[
 j(\operatorname{Alb}S_1)=J(-w^2/48^3).
\]
This function is holomorphic on $0<|w|<48^{3/2}$, has a pole of order
two at zero, and has a natural boundary at $|w|=48^{3/2}$.
In particular it continues holomorphically through every point of the
family boundary $|w|=216$.
\end{statement}
\begin{proof}
The precise Albanese lattices in R5 are verified by projecting all
translations on the product covers. At the order-three fibre the
projection is $(1/3)(\mathbb Z+\kappa_1\mathbb Z)$; multiplication by
three gives the first stated lattice. At the order-four fibre it is
$\tfrac12\mathbb Z+\tfrac14\kappa_2\mathbb Z$; multiplication by two
gives $\mathbb Z+(\kappa_2/2)\mathbb Z$. Thus the positive-imaginary
period representatives are $-\kappa_1$ and $-\kappa_2/2$.
Adding $2\mu_1=1-i/\sqrt3$ and
$3\mu_2=3/2-3i/2$ to C.1.2 proves the two $\kappa$ identities,
including their real parts. Exponentiation gives the signs in $q_1,q_2$;
the ratio is $216^2/48^3=27/64$.

We checked checkpoint K2's step giving $q_2=iw/216$: its factor
$1/2$ is exactly the Albanese projection just computed. Its
natural-boundary proof uses modular invariance at rational boundary
points and is valid. The existence criterion and exact $2\mathbb Z$
identification are R4 and R5; the latter is also Theorem~1.4
(label \texttt{p2:thm:exact-period}) of
\cite{DuGuoKusakabeXie2026v1}. Its hypotheses specify precisely the
retained period family, affine finite twists and cusp collar. They
coincide with the family in this prompt; the normalization differs
only by the fixed choice of the particular solution, which does not
alter parameter differences. Version 2 drops that theorem; we infer
no mathematical conclusion from this omission.

The Fourier expansion for $J$ converges for $0<|q|<1$ and has a simple
pole at $q=0$. The inverse image under $q=-w^2/48^3$ is the claimed
punctured disc and the pole has order two. The inequality
$216<48^{3/2}$ is equivalent to $27<64$. At every nonzero point on the
larger boundary the map to $q$ is locally biholomorphic. A continuation
there would continue $J$ through the unit circle, contradicting the
same rational-cusp argument checked in K2. The phrase ``up to the
radius'' therefore means the open punctured disc, not holomorphy on
its bounding circle or at its central pole.
\end{proof}
